\begin{definition}
    \textbf{Successione regolare}

    Una successione $\left\{x_n\right\}_{n\in\mathbb{N}}$ in uno spazio
    $\left(\mathbb{R}, \left|\cdot\right|\right)$ si dice regolare se è convergente
    oppure divergente.
\end{definition}

\begin{proposition}
    Sia $\left\{x_n\right\}_{n\in\mathbb{N}}$ una successione regolare,
    allora il limite è unico ed e anche limite di ogni sua
    sottosuccessione. In simboli,
    \[
        \underset{n\rightarrow+\infty}{\lim} x_n = l\in\left[-\infty,+\infty\right]
        \Rightarrow \begin{cases}
            l \text{ è unico}\\
            \forall \left\{x_{n_k}\right\}_k \Rightarrow \underset{k\rightarrow+\infty}{\lim}x_{n_k} = l
        \end{cases}
    \]
\end{proposition}

\begin{observation}
    Una successione divergente non è di Cacuchy, poiché non è convergente.
\end{observation}

\begin{observation}
    La successione $\left\{x_n\right\}_{n\in\mathbb{N}}, x_n := \left(-1\right)^n n$ NON è
    divergente. Difatti due ha due diverse sottosuccessioni che tnedono a valori diversi,
    se effettivamente fosse divergente il valore del limite non potrebbe risualtare unico
    se deve essere anche limite di tutte le sottosuccessioni.

    Difatti
    \[
        n_k = 2k  \implies x_{n_k} = \left(-1\right)^{2k} 2k = \left(\left(-1\right)^2\right)^k 2k =
        \left(1\right) ^ k 2k = 1 \cdot 2k = 2k \implies\]\[
        \implies \underset{k\rightarrow+\infty}{\lim} x_{n_k} = +\infty,
    \]\[
        n_k = 2k+1  \implies x_{n_k} = \left(-1\right)^{2k+1} \left(2k+1\right) =
        \left(-1\right)^{2k} \left(-1\right)^{1} \left(2k+1\right) = \]\[
        = \left(\left(-1\right)^2\right)^k \left(-1\right) \left(2k+1\right) =
        - \left(1\right) ^ k \left(2k+1\right) =
        - \left(1\right) \left(2k+1\right) = \]\[
        = -\left(2k+1\right) \implies \underset{k\rightarrow+\infty}{\lim} x_{n_k} = -\infty.
    \]
\end{observation}